% CHECK: Claude's translation of the English abstract; please read it as a native speaker before submitting.
Neuronové sítě, které čtou přímo obrázky, dokážou řešit generované Ravenovy progresivní matice, ale jen s~velkým
množstvím dat: vision transformer trénovaný na I-RAVEN vyřeší 19\,\% úloh při 1\,000 trénovacích úlohách na každé
rozvržení a 93\,\% při 6\,000. Tato práce studuje alternativu, tedy strukturu vestavěnou do modelu před trénováním, a
ptá se, jakou cenu má pro relační model vestavěná symetrie a zda ji lze poznat ještě před trénováním.

Modely jsou liftované relační neuronové sítě v~PyNeuraLogic, v~nichž se šablona vážených pravidel pro každou úlohu
zakotví (grounding) do jiné sítě. Dokazujeme, že šablona nedokáže rozlišit konstanty, které nepojmenovává. Symetrii
lze tedy vestavět tím, že transformované věci zůstanou nepojmenované, a pak platí přesně pro všechna nastavení vah.
Touto metodou stavíme model, který generuje chybějící panel z~grafu objektů úlohy, aniž by dostal pravidla nebo
viděl kandidátní odpovědi. Na testovací sadě I-RAVEN dosahuje 97,2\,\% z~popisů objektů a 97,0\,\% end-to-end
s~rozpoznávačem objektů, který se učí z~označených objektů každého panelu. Vision transformer, který dostává pouze
označení odpovědi, dosahuje 93,1\,\% se šestinásobkem vah.

Poté symetriím přiřazujeme cenu. Spočítání orbit tabulky pravidla pod symetrií modelu předpovídá, kolik dat
symetrie ušetří a jaký strop nastaví, když ji pravidlo nesdílí. Ověřují to dvě pravidla shodná formou i velikostí,
součet a střed. Z~600 úloh dosahuje model ekvivariantní vůči posunu, který čte odpověď z~řádku s~otázkou,
91\,\% na středu a 6\,\% na součtu, jehož předpovězený strop je 7,9\,\%; když dostane sčítání jako znalost na pozadí, naučí se součet ze 100 úloh. Na
benchmarku stojí odebrání symetrie posunu 61 procentních bodů přesnosti při 100 trénovacích úlohách a 2 body při
6\,000. Symetrie mezi atributy objektu má, jak bylo spočteno před během experimentů, cenu dvoj- až trojnásobku
trénovacích dat, a model bez ní jich tolik potřebuje navíc. Symetrie, která je v~rozporu s~daty, drží model u~stropu
spočteného předem. Počty omezují, čeho může síť dosáhnout; zda toho dosáhne, závisí i~na tom, jak úlohu čte, a práce
uvádí, kde za nimi zaostává.
